\chapter{Les catégories de complexes des catégories additives}
\VerdierPhysicalPageAnchor{0034}
\VerdierSemanticAnchor{ch1}

\VerdierSemanticAnchor{ch1-1}
\section{Catégories graduées}

\VerdierSemanticAnchor{ch1-1-1}
\subsection{Catégories graduées de type \texorpdfstring{$G$}{G}}

\VerdierSemanticAnchor{ch1-1-1-1}
\noindent\textbf{1.1.1.}\quad
À tout groupe $G$ on associe une catégorie, que nous notons $\mathcal G$,
définie de la manière suivante :

\VerdierSemanticAnchor{ch1-1-1-1-group-category}
\noindent\textnormal{(1.1.1.1)}
\begin{list}{--}{%
  \setlength{\leftmargin}{2em}%
  \setlength{\itemsep}{0.25em}%
  \setlength{\topsep}{0.25em}%
}
  \item \VerdierSemanticAnchor{ch1-1-1-1-group-category-ob}%
    $\operatorname{Ob}(\mathcal G)$ est un ensemble réduit à un élément
    choisi une fois pour toutes.
  \item \VerdierSemanticAnchor{ch1-1-1-1-group-category-fl}%
    $\operatorname{Fl}(\mathcal G)$ est l'ensemble sous-jacent à $G$.
  \item \VerdierSemanticAnchor{ch1-1-1-1-group-category-maps}%
    Les applications source et but sont les uniques applications de
    $\operatorname{Fl}(\mathcal G)$ dans $\operatorname{Ob}(\mathcal G)$.
    L'application identité envoie l'unique élément de
    $\operatorname{Ob}(\mathcal G)$ sur l'élément neutre de $G$.
  \item \VerdierSemanticAnchor{ch1-1-1-1-group-category-composition}%
    Le composé de deux flèches
    ${}\xrightarrow{g}{}\xrightarrow{h}{}$ est la flèche $hg$, produit dans
    le groupe $G$ des éléments $h$ et $g$.
\end{list}

\VerdierSemanticAnchor{ch1-1-1-1-graded-category}
On appelle \emph{catégorie graduée de type $G$} une catégorie fibrée sur
$\mathcal G$ [18]. Une catégorie $(\mathcal C,G)$ graduée de type $G$ est
donc une catégorie $\mathcal C$ munie d'un foncteur
$\deg\colon\mathcal C\to\mathcal G$ faisant de $\mathcal C$ une catégorie
fibrée au-dessus de $\mathcal G$.

\VerdierSemanticAnchor{ch1-1-1-1-degree-functor}
Le foncteur $\deg\colon\mathcal C\to\mathcal G$ est appelé
\emph{foncteur degré}. L'image par $\deg$ d'un morphisme de $\mathcal C$ est
appelé le \emph{degré de ce morphisme}.
\VerdierSemanticAnchor{ch1-1-1-1-fiber-category}
Soit $e$ l'élément neutre de $G$.
On désigne par $\mathcal C_e$ la catégorie fibre de $(\mathcal C,G)$. La
catégorie $\mathcal C_e$ est
